\documentclass[10pt]{article}
\usepackage[letterpaper,margin=0.72in]{geometry}
\usepackage[T1]{fontenc}
\usepackage{lmodern,microtype,xcolor,enumitem,hyperref}
\definecolor{navy}{HTML}{08111F}
\definecolor{blue}{HTML}{2677FF}
\definecolor{muted}{HTML}{5D6B80}
\hypersetup{colorlinks=true,urlcolor=blue}
\setlist[itemize]{leftmargin=1.2em,itemsep=2pt,topsep=3pt}
\pagestyle{plain}
\begin{document}
{\color{blue}\small\bfseries SUMMER ENGINEERING PROJECT --- 2026}
\vspace{0.4em}
{\color{navy}\LARGE\bfseries Ramjet and Scramjet Simulation Concepts}\\
{\color{muted}\large One-dimensional propulsion-cycle models and design-space exploration}

\section*{Project overview}
A conceptual simulation study for air-breathing high-speed propulsion. The project organizes inlet compression, heat addition, combustor constraints, and nozzle expansion into traceable one-dimensional models for ramjet and scramjet regimes.

\section*{Engineering problem}
High-speed propulsion performance depends on tightly coupled compressible-flow processes, and ideal-cycle results can be misleading unless losses, thermal limits, and operating assumptions are explicit.

\section*{System architecture}
\begin{itemize}
\item Atmosphere and flight condition define freestream stagnation properties.
\item Inlet/diffuser model applies compression with configurable pressure recovery.
\item Combustor model applies heat addition, fuel limits, and ramjet/scramjet Mach assumptions.
\item Nozzle model estimates expansion, exit state, thrust, and specific impulse trends.
\end{itemize}

\section*{Engineering focus}
\begin{itemize}
\item Track static and stagnation properties station by station with unit checks.
\item Compare ideal and loss-inclusive cases across Mach number and equivalence ratio.
\item Treat dissociation, finite-rate chemistry, and multidimensional shocks as future fidelity.
\end{itemize}

\section*{Verification plan}
\begin{itemize}
\item Recover limiting textbook cases before exploring the design space.
\item Check conservation of mass, energy, and physically admissible temperatures and pressures.
\item Use sensitivity plots to distinguish robust trends from assumption-driven conclusions.
\end{itemize}

\section*{Status and evidence boundary}
Summer 2026 simulation concept; this brief presents model structure, not validated engine performance.

\vfill
\small\color{muted} V. Narayan Kushwaha --- Summer Engineering Portfolio 2026
\end{document}
